\section{The Owner as Transformer: Lessons from Bending Spoons}\label{sec:acquirer}
% Generated by code/verify_part2_acquirer.py from notes/case_bending_spoons_f1_extract.md and data/examples/invoice_chain.json (version 2026-10-29c); do not edit.
\providecommand{\acqAcqHOneSix}{\$2{,}286 million}
\providecommand{\acqAcqShareHOneFive}{99.9\%}
\providecommand{\acqAcqShareHOneSix}{99.8\%}
\providecommand{\acqAcqShareTwentyFive}{100.0\%}
\providecommand{\acqAcqShareTwentyFiveTwo}{99.97\%}
\providecommand{\acqAcqShareTwentyFour}{81.7\%}
\providecommand{\acqAirEq}{\$2.25 billion}
\providecommand{\acqAirEv}{\$1.285 billion}
\providecommand{\acqAmzQOne}{11.5\%}
\providecommand{\acqAmzTwentyFive}{11.5\%}
\providecommand{\acqAmzTwentyFour}{12.8\%}
\providecommand{\acqAnnFactor}{0.51}
\providecommand{\acqAnnPtsAnnLo}{0.10}
\providecommand{\acqAnnPtsHi}{0.47}
\providecommand{\acqAnnPtsLo}{0.19}
\providecommand{\acqAoimQOne}{51.2\%}
\providecommand{\acqAoimTwentyFive}{47\%}
\providecommand{\acqAoimTwentyFiveX}{46.9\%}
\providecommand{\acqAoimTwentyFour}{44.6\%}
\providecommand{\acqAolB}{\$1.5 billion}
\providecommand{\acqAolDebt}{\$2.8 billion}
\providecommand{\acqAolEmp}{214}
\providecommand{\acqAolOpm}{53\%}
\providecommand{\acqAolWarn}{108}
\providecommand{\acqApps}{800{,}000}
\providecommand{\acqBaseTwentyFive}{5\%}
\providecommand{\acqBcCut}{198}
\providecommand{\acqBcPrice}{\$233 million}
\providecommand{\acqBcPs}{1.2}
\providecommand{\acqBcPx}{\$4.45}
\providecommand{\acqBcStaff}{600}
\providecommand{\acqBoundFac}{0.51}
\providecommand{\acqBrkOne}{\$600 million}
\providecommand{\acqBrkStake}{4\%}
\providecommand{\acqBrkTwo}{\$330 million}
\providecommand{\acqCapClose}{\$25.7 billion}
\providecommand{\acqCapOpen}{\$18.4 billion}
\providecommand{\acqCapexHOneSix}{\$4.0 million}
\providecommand{\acqCapexRevTwentyFive}{0.04\%}
\providecommand{\acqCapexRevTwentyFour}{0.09\%}
\providecommand{\acqCapexShareHOneFive}{0.05\%}
\providecommand{\acqCapexShareHOneSix}{0.2\%}
\providecommand{\acqCffShare}{148\%}
\providecommand{\acqCffTwentyFive}{\$1{,}935.9 million}
\providecommand{\acqCfiAbsTwentyFive}{\$1{,}851.9 million}
\providecommand{\acqCfiQOne}{274.1\%}
\providecommand{\acqCfiTwentyFive}{141.8\%}
\providecommand{\acqCfiTwentyFour}{132.6\%}
\providecommand{\acqChCouponExpHi}{113.7}
\providecommand{\acqChCouponExpLo}{86.7}
\providecommand{\acqChCouponHi}{112.8}
\providecommand{\acqChCouponLo}{86.0}
\providecommand{\acqChPaidHi}{113.1}
\providecommand{\acqChPaidLo}{86.6}
\providecommand{\acqConvHi}{114}
\providecommand{\acqConvLevHi}{103}
\providecommand{\acqConvLevLo}{79}
\providecommand{\acqConvLo}{87}
\providecommand{\acqCouponEur}{8.625}
\providecommand{\acqCouponMarg}{9.625}
\providecommand{\acqCov}{4.00}
\providecommand{\acqCrfHi}{10.8\%}
\providecommand{\acqCrfLo}{14.6\%}
\providecommand{\acqCumOutTwentyFive}{2.23}
\providecommand{\acqCumOutTwentyFour}{1.58}
\providecommand{\acqDayOne}{39.7\%}
\providecommand{\acqDeclHi}{9\%}
\providecommand{\acqDeclLo}{5\%}
\providecommand{\acqDepPrevQOne}{1.5}
\providecommand{\acqDepQOne}{\$2.01 billion}
\providecommand{\acqDepRatio}{10}
\providecommand{\acqDepTwentyThree}{\$194 million}
\providecommand{\acqDilution}{6.1\%}
\providecommand{\acqDoubleTwentyFive}{8.3}
\providecommand{\acqDoubleZeroTwentyFive}{10.9}
\providecommand{\acqEbM}{\$500 million}
\providecommand{\acqEbPx}{\$4.50}
\providecommand{\acqEqRaised}{\$549 million}
\providecommand{\acqEqShare}{99\%}
\providecommand{\acqEquityBs}{\$1.06 billion}
\providecommand{\acqFlagCash}{25\%}
\providecommand{\acqFlagCorr}{17}
\providecommand{\acqFlagCov}{26\%}
\providecommand{\acqFlagCovFast}{41\%}
\providecommand{\acqFlagG}{50\%}
\providecommand{\acqFlagPf}{10\%}
\providecommand{\acqFteTotal}{2{,}300}
\providecommand{\acqFullAnnAol}{\ensuremath{-}0.18}
\providecommand{\acqFullAnnEb}{\ensuremath{+}0.11}
\providecommand{\acqFullAnnVim}{\ensuremath{-}0.66}
\providecommand{\acqFullPtsHi}{2.2}
\providecommand{\acqFullPtsLo}{0.7}
\providecommand{\acqGeoEu}{21\%}
\providecommand{\acqGeoNa}{65\%}
\providecommand{\acqGiAol}{\ensuremath{-}5.1\%}
\providecommand{\acqGiEvernote}{\ensuremath{-}1.0\%}
\providecommand{\acqGiRemini}{\ensuremath{-}13.9\%}
\providecommand{\acqGiStreamyard}{\ensuremath{-}9.4\%}
\providecommand{\acqGiTwentyFive}{\ensuremath{-}5.1\%}
\providecommand{\acqGiTwentyFour}{\ensuremath{-}9.4\%}
\providecommand{\acqGinHi}{0.10}
\providecommand{\acqGinLo}{0.03}
\providecommand{\acqGmaxHi}{12\%}
\providecommand{\acqGmaxLo}{1\%}
\providecommand{\acqGmaxOrgHi}{20\%}
\providecommand{\acqGmaxOrgLo}{9\%}
\providecommand{\acqGridLevMax}{42\%}
\providecommand{\acqGridLevMin}{34\%}
\providecommand{\acqGridMax}{1.82}
\providecommand{\acqGridMfiveA}{1.82}
\providecommand{\acqGridMfourA}{1.46}
\providecommand{\acqGridMfourB}{1.11}
\providecommand{\acqGridMthreeA}{1.09}
\providecommand{\acqGridPfourA}{1.11}
\providecommand{\acqGridPfourB}{0.92}
\providecommand{\acqGrowthQOne}{132.2\%}
\providecommand{\acqGrowthTwentyFive}{94.7\%}
\providecommand{\acqGrowthTwentyFour}{73.4\%}
\providecommand{\acqHarvestMult}{15.2}
\providecommand{\acqHarvestNew}{\$2{,}110}
\providecommand{\acqHarvestOld}{\$130}
\providecommand{\acqHarvestRise}{1{,}523\%}
\providecommand{\acqHazFive}{0.147}
\providecommand{\acqHazMean}{0.125}
\providecommand{\acqHazTen}{0.127}
\providecommand{\acqHireRate}{0.036\%}
\providecommand{\acqHires}{286}
\providecommand{\acqHurL}{65\%}
\providecommand{\acqHurU}{25\%}
\providecommand{\acqIcCpool}{\$185.0 thousand}
\providecommand{\acqIcCrun}{\$329.0 thousand}
\providecommand{\acqIcCspec}{\$144.0 thousand}
\providecommand{\acqIcF}{\$2.81 million}
\providecommand{\acqIcFloorFour}{$-\$570$ thousand}
\providecommand{\acqIcFloorFourCertain}{$-\$458$ thousand}
\providecommand{\acqIcFloorInf}{\$141 thousand}
\providecommand{\acqIcFloorVone}{\$408 thousand}
\providecommand{\acqIcFreuse}{\$2.62 million}
\providecommand{\acqIcLimNoReb}{\$20.9 thousand}
\providecommand{\acqIcNbe}{7}
\providecommand{\acqIcNmin}{2.1}
\providecommand{\acqIcNminNoReb}{11.2}
\providecommand{\acqIcNpv}{$-\$3.78$ million}
\providecommand{\acqIcTh}{4}
\providecommand{\acqIcThbe}{12}
\providecommand{\acqIcUFbe}{0.84}
\providecommand{\acqIcVarpi}{0.82}
\providecommand{\acqIcVarpiBe}{0.73}
\providecommand{\acqIcW}{$-\$164.1$ thousand}
\providecommand{\acqIcWzero}{\$164.9 thousand}
\providecommand{\acqIntAbsQOne}{\$93.2 million}
\providecommand{\acqIntAbsTwentyFive}{\$142.6 million}
\providecommand{\acqIntPaidHOneSix}{\$161.2 million}
\providecommand{\acqIntPaidRatioHOneFive}{90\%}
\providecommand{\acqIntPaidRatioHOneSix}{80\%}
\providecommand{\acqIntPaidShQOne}{14.8\%}
\providecommand{\acqIntPaidShTwentyFive}{10.0\%}
\providecommand{\acqIntPaidShTwentyFour}{2.9\%}
\providecommand{\acqIntPaidShTwentyThree}{4.3\%}
\providecommand{\acqIntPaidTwentyFive}{\$130.9 million}
\providecommand{\acqIntQOne}{15.5\%}
\providecommand{\acqIntTwentyFive}{10.9\%}
\providecommand{\acqIntTwentyFour}{4.9\%}
\providecommand{\acqIntangShareTwentyFour}{18.3\%}
\providecommand{\acqIntangTwentyFour}{\$162.5 million}
\providecommand{\acqIpoCompany}{\$1.0 billion}
\providecommand{\acqIpoCompanyShares}{34.4 million}
\providecommand{\acqIpoNet}{\$1.10 billion}
\providecommand{\acqIpoNetEst}{\$933 million}
\providecommand{\acqIpoNetFull}{\$1{,}079 million}
\providecommand{\acqIpoNewP}{34.40 million}
\providecommand{\acqIpoPx}{\$29.00}
\providecommand{\acqIpoShares}{57.97 million}
\providecommand{\acqIpoSharesP}{57.97 million}
\providecommand{\acqIpoTotal}{\$1.68 billion}
\providecommand{\acqIrrL}{50.1\%}
\providecommand{\acqIrrPe}{19.9\%}
\providecommand{\acqIrrU}{22.3\%}
\providecommand{\acqKAol}{2.3}
\providecommand{\acqKEb}{1.7}
\providecommand{\acqKHi}{3.2}
\providecommand{\acqKLo}{1.7}
\providecommand{\acqKVim}{3.2}
\providecommand{\acqKcapHi}{2.1}
\providecommand{\acqKcapLo}{1.2}
\providecommand{\acqKomootMax}{85\%}
\providecommand{\acqKomootStaff}{150}
\providecommand{\acqKucoin}{\$13.1 billion}
\providecommand{\acqLevQOne}{2.19}
\providecommand{\acqLevQTwo}{2.4}
\providecommand{\acqLevTwentyFive}{2.24}
\providecommand{\acqLiabShare}{85\%}
\providecommand{\acqLifeHi}{19.5}
\providecommand{\acqLifeLo}{10.6}
\providecommand{\acqLogGrowthTwentyFive}{66.6\%}
\providecommand{\acqLtv}{85\%}
\providecommand{\acqMaiQOne}{4.5\%}
\providecommand{\acqMaiTwentyFive}{10.4\%}
\providecommand{\acqMaiTwentyFour}{14.1\%}
\providecommand{\acqMaxConcurrent}{20}
\providecommand{\acqMaxConcurrentExact}{20.7}
\providecommand{\acqMcfUQOne}{24.4\%}
\providecommand{\acqMcfUTwentyFive}{30.5\%}
\providecommand{\acqMcfUTwentyFour}{34.2\%}
\providecommand{\acqMgridHi}{0.3}
\providecommand{\acqMgridLo}{0.2}
\providecommand{\acqMiroEq}{\$1.79 billion}
\providecommand{\acqMiroEv}{\$1.355 billion}
\providecommand{\acqMiroReinv}{\$295 million}
\providecommand{\acqMixFast}{0.54}
\providecommand{\acqMixNrr}{83\%}
\providecommand{\acqMixSlow}{0.10}
\providecommand{\acqMixW}{22\%}
\providecommand{\acqMosaic}{330}
\providecommand{\acqMwBunch}{9.2\%}
\providecommand{\acqMwBunchLate}{31.4\%}
\providecommand{\acqMwCross}{2.46\%}
\providecommand{\acqMwCrossRise}{4\%}
\providecommand{\acqMwLate}{9.0\%}
\providecommand{\acqMwOne}{2.4\%}
\providecommand{\acqMwOneExact}{2.37\%}
\providecommand{\acqMwSpread}{9.5\%}
\providecommand{\acqMwZ}{0.176}
\providecommand{\acqMwZdays}{3}
\providecommand{\acqNeedTwentyFive}{9\%}
\providecommand{\acqNeedTwentyFour}{11\%}
\providecommand{\acqNeededTwentyFive}{99.7\%}
\providecommand{\acqNetDebt}{\$4.09 billion}
\providecommand{\acqNiTwentyFive}{\$\ensuremath{-}0.2 million}
\providecommand{\acqNotesCoupon}{168}
\providecommand{\acqNotesDue}{2032}
\providecommand{\acqNotesEur}{EUR 500 million}
\providecommand{\acqNotesEurC}{8.625\%}
\providecommand{\acqNotesUsd}{\$1.25 billion}
\providecommand{\acqNotesUsdC}{9.625\%}
\providecommand{\acqNrrAol}{95\%}
\providecommand{\acqNrrEvernote}{99\%}
\providecommand{\acqNrrHi}{95\%}
\providecommand{\acqNrrLo}{91\%}
\providecommand{\acqNrrQOne}{94\%}
\providecommand{\acqNrrRemini}{87\%}
\providecommand{\acqNrrStreamyard}{91\%}
\providecommand{\acqNrrTwentyFive}{95\%}
\providecommand{\acqNrrTwentyFour}{91\%}
\providecommand{\acqNrrTwentyThree}{93\%}
\providecommand{\acqNuTwentyFive}{0.051}
\providecommand{\acqNuTwentyFour}{0.094}
\providecommand{\acqOcfCagr}{122\%}
\providecommand{\acqOcfCfiQOne}{\ensuremath{-}261.5\%}
\providecommand{\acqOcfCfiTwentyFive}{\ensuremath{-}119.5\%}
\providecommand{\acqOcfCfiTwentyFour}{\ensuremath{-}102.0\%}
\providecommand{\acqOcfGrowth}{246\%, 42\% and 254\%}
\providecommand{\acqOcfTwentyFive}{\$290.6 million}
\providecommand{\acqOcfmQOne}{12.6\%}
\providecommand{\acqOcfmTwentyFive}{22.2\%}
\providecommand{\acqOcfmTwentyFour}{30.5\%}
\providecommand{\acqOpmQOne}{20.0\%}
\providecommand{\acqOpmTwentyFive}{21.3\%}
\providecommand{\acqOpmTwentyFour}{19.0\%}
\providecommand{\acqOrgGTwentyFive}{13\%}
\providecommand{\acqOrgGTwentyFour}{7\%}
\providecommand{\acqOrgLogQTwo}{3.0\%}
\providecommand{\acqOrgPage}{80}
\providecommand{\acqOrgQOne}{83\%}
\providecommand{\acqOrgTwentyFive}{79\%}
\providecommand{\acqOrgTwentyFour}{76\%}
\providecommand{\acqOutPrevTwentyFive}{2.8}
\providecommand{\acqOutPrevTwentyFour}{2.3}
\providecommand{\acqPfExcess}{100\%}
\providecommand{\acqPfInt}{14.1\%}
\providecommand{\acqPfOpm}{19.5\%}
\providecommand{\acqPfPage}{143}
\providecommand{\acqPfRatio}{2.0}
\providecommand{\acqPfRev}{\$2{,}608.5 million}
\providecommand{\acqPninetyfiveHi}{13}
\providecommand{\acqPninetyfiveLo}{7}
\providecommand{\acqPrHi}{90\%}
\providecommand{\acqPrLo}{10\%}
\providecommand{\acqPrSolo}{70\%}
\providecommand{\acqPrivRound}{\$11 billion}
\providecommand{\acqPxAol}{\$1{,}454 million}
\providecommand{\acqPxEb}{\$505 million}
\providecommand{\acqPxVim}{\$1{,}359 million}
\providecommand{\acqQTwoAdjM}{54\%}
\providecommand{\acqQTwoAdjMx}{54.1\%}
\providecommand{\acqQTwoAmz}{11.7\%}
\providecommand{\acqQTwoG}{126\%}
\providecommand{\acqQTwoGx}{126.3\%}
\providecommand{\acqQTwoInt}{15.5\%}
\providecommand{\acqQTwoIntAbs}{\$109.0 million}
\providecommand{\acqQTwoM}{34\%}
\providecommand{\acqQTwoMx}{34.1\%}
\providecommand{\acqQTwoOrg}{3\%}
\providecommand{\acqQTwoReorg}{\$51 million}
\providecommand{\acqQTwoReorgAbs}{\$50.8 million}
\providecommand{\acqQTwoReorgSh}{7.2\%}
\providecommand{\acqQTwoReorgX}{7.2\%}
\providecommand{\acqQTwoRev}{\$704.2 million}
\providecommand{\acqQTwoTrx}{0.9\%}
\providecommand{\acqRAdjQOne}{183.4}
\providecommand{\acqRAdjTwentyFive}{141.6}
\providecommand{\acqRAdjTwentyFour}{118.0}
\providecommand{\acqRAllInvQOne}{\ensuremath{-}129.3}
\providecommand{\acqRAllInvTwentyFive}{\ensuremath{-}24.8}
\providecommand{\acqRAllInvTwentyFour}{\ensuremath{-}28.7}
\providecommand{\acqRAllInvUQOne}{\ensuremath{-}117.5}
\providecommand{\acqRAllInvUTwentyFive}{\ensuremath{-}16.5}
\providecommand{\acqRAllInvUTwentyFour}{\ensuremath{-}25.0}
\providecommand{\acqRAnnHiTwentyFive}{92.8}
\providecommand{\acqRAnnHiTwentyFour}{86.8}
\providecommand{\acqRAnnLoTwentyFive}{84.4}
\providecommand{\acqRAnnLoTwentyFour}{80.9}
\providecommand{\acqRGaapQOne}{152.2}
\providecommand{\acqRGaapTwentyFive}{115.9}
\providecommand{\acqRGaapTwentyFour}{92.3}
\providecommand{\acqRIntQOne}{136.7}
\providecommand{\acqRIntTwentyFive}{105.0}
\providecommand{\acqRIntTwentyFour}{87.5}
\providecommand{\acqRLaAvgHiTwentyFive}{111.3}
\providecommand{\acqRLaAvgHiTwentyFour}{97.7}
\providecommand{\acqRLaAvgLoTwentyFive}{104.4}
\providecommand{\acqRLaAvgLoTwentyFour}{94.2}
\providecommand{\acqRLaFullHiTwentyFive}{101.1}
\providecommand{\acqRLaFullHiTwentyFour}{90.5}
\providecommand{\acqRLaFullLoTwentyFive}{92.7}
\providecommand{\acqRLaFullLoTwentyFour}{84.5}
\providecommand{\acqROcfQOne}{144.8}
\providecommand{\acqROcfSbcOnlyTwentyFive}{112.7}
\providecommand{\acqROcfSbcTwentyFive}{112.6}
\providecommand{\acqROcfSbcTwentyFour}{100.0}
\providecommand{\acqROcfTwentyFive}{116.9}
\providecommand{\acqROcfTwentyFour}{103.9}
\providecommand{\acqROcfUQOne}{156.6}
\providecommand{\acqROcfUTwentyFive}{125.2}
\providecommand{\acqROcfUTwentyFour}{107.6}
\providecommand{\acqROcfUZeroTwentyFive}{127.8}
\providecommand{\acqROrgIntNoReorgTwentyFive}{29.4}
\providecommand{\acqROrgIntNoReorgTwentyFour}{28.8}
\providecommand{\acqROrgIntTwentyFive}{23.4}
\providecommand{\acqROrgIntTwentyFiveR}{23}
\providecommand{\acqROrgIntTwentyFour}{21.1}
\providecommand{\acqROrgNoReorgTwentyFive}{40.3}
\providecommand{\acqROrgNoReorgTwentyFiveR}{40}
\providecommand{\acqROrgNoReorgTwentyFour}{33.7}
\providecommand{\acqROrgTwentyFive}{34.3}
\providecommand{\acqROrgTwentyFour}{26.0}
\providecommand{\acqRQTwoAdj}{180.5}
\providecommand{\acqRQTwoGap}{123}
\providecommand{\acqRQTwoOrg}{37}
\providecommand{\acqRQTwoOrgInt}{21.6}
\providecommand{\acqRQTwoOrgIntNoReorg}{28.9}
\providecommand{\acqRQTwoOrgIntNoReorgR}{29}
\providecommand{\acqRQTwoOrgIntR}{22}
\providecommand{\acqRQTwoOrgNoReorg}{44.3}
\providecommand{\acqRQTwoOrgNoReorgR}{44}
\providecommand{\acqRQTwoOrgX}{37.1}
\providecommand{\acqRQTwoRep}{160}
\providecommand{\acqRQTwoRepInt}{145.0}
\providecommand{\acqRQTwoRepIntR}{145}
\providecommand{\acqRQTwoRepX}{160.5}
\providecommand{\acqRUcAvgHiTwentyFive}{108.0}
\providecommand{\acqRUcAvgHiTwentyFour}{95.0}
\providecommand{\acqRUcAvgLoTwentyFive}{100.9}
\providecommand{\acqRUcAvgLoTwentyFour}{91.5}
\providecommand{\acqRUcEndHiTwentyFive}{95.2}
\providecommand{\acqRUcEndHiTwentyFour}{86.1}
\providecommand{\acqRUcEndLoTwentyFive}{87.2}
\providecommand{\acqRUcEndLoTwentyFour}{80.3}
\providecommand{\acqRdebt}{9.19\%}
\providecommand{\acqRdebtReg}{7.7\%}
\providecommand{\acqReminiFte}{30}
\providecommand{\acqReorgAbsTwentyFive}{\$78.6 million}
\providecommand{\acqReorgQOne}{12.6\%}
\providecommand{\acqReorgRatio}{3.6}
\providecommand{\acqReorgTwentyFive}{6.0\%}
\providecommand{\acqReorgTwentyFour}{7.7\%}
\providecommand{\acqReorgUsdEx}{\$30.1 million}
\providecommand{\acqReqCashHi}{246\%}
\providecommand{\acqReqCashLo}{124\%}
\providecommand{\acqReqHi}{76.0\%}
\providecommand{\acqReqLo}{71.7\%}
\providecommand{\acqReqOrg}{63.7\%}
\providecommand{\acqRetHi}{28\%}
\providecommand{\acqRetLo}{22\%}
\providecommand{\acqRetNetHi}{22\%}
\providecommand{\acqRetNetLo}{14\%}
\providecommand{\acqRetNetTauHi}{19.8\%}
\providecommand{\acqRetNetTauLo}{13.6\%}
\providecommand{\acqRetTauHi}{25.8\%}
\providecommand{\acqRetTauLo}{22.3\%}
\providecommand{\acqRevQOneFive}{\$258.9 million}
\providecommand{\acqRevQOneSix}{\$601.3 million}
\providecommand{\acqRevTwentyFive}{\$1{,}306.4 million}
\providecommand{\acqRevTwentyFour}{\$671.1 million}
\providecommand{\acqRevTwentyThree}{\$387.1 million}
\providecommand{\acqRpCustEight}{0.21\%}
\providecommand{\acqRpCustOne}{1.69\%}
\providecommand{\acqRpCustXbOne}{0.67\%}
\providecommand{\acqRpCustXbThree}{0.19\%}
\providecommand{\acqRpCustXbTwo}{0.34\%}
\providecommand{\acqRpFteDE}{387}
\providecommand{\acqRpFteFR}{705}
\providecommand{\acqRpFteUK}{1{,}124}
\providecommand{\acqRpFteUS}{337}
\providecommand{\acqRpGainBound}{11.40}
\providecommand{\acqRpGainDE}{6.45}
\providecommand{\acqRpGainFR}{5.87}
\providecommand{\acqRpGainHalf}{5.53}
\providecommand{\acqRpGainUS}{8.43}
\providecommand{\acqRpsQOne}{\$0.97 million}
\providecommand{\acqRpsTwentyFive}{\$2.57 million}
\providecommand{\acqRpsTwentyThree}{\$1.12 million}
\providecommand{\acqRtv}{8.93\%}
\providecommand{\acqRunRateInt}{572}
\providecommand{\acqRunRateIntHOne}{404}
\providecommand{\acqRunRateIntHi}{787}
\providecommand{\acqRunRateShPf}{14.7\%}
\providecommand{\acqRunRateShPfHi}{20.2\%}
\providecommand{\acqRunRateShQ}{20.3\%}
\providecommand{\acqRunRateShQHi}{27.9\%}
\providecommand{\acqRw}{8.93\%}
\providecommand{\acqRwAnnual}{9.34\%}
\providecommand{\acqSbcTwentyFive}{4.3\%}
\providecommand{\acqSbcTwentyFiveTwo}{4.25\%}
\providecommand{\acqSbcTwentyFour}{3.9\%}
\providecommand{\acqSbcTwentyThree}{3.8\%}
\providecommand{\acqSbcUsdTwentyFive}{\$55.5 million}
\providecommand{\acqScoopLoss}{\$84.5 million}
\providecommand{\acqSpo}{621}
\providecommand{\acqSpoAvgTwentyFive}{508}
\providecommand{\acqSpoAvgTwentyThree}{346}
\providecommand{\acqSpoShare}{27\%}
\providecommand{\acqSpoTech}{42\%}
\providecommand{\acqStandCashHi}{30.5\%}
\providecommand{\acqStandCashLo}{8.9\%}
\providecommand{\acqStandHi}{5.1\%}
\providecommand{\acqStandLo}{9.4\%}
\providecommand{\acqStandUsdHi}{\$67 million}
\providecommand{\acqStandUsdLo}{\$123 million}
\providecommand{\acqSubsTwentyFive}{93\%}
\providecommand{\acqSubsTwentyFour}{92\%}
\providecommand{\acqSubsTwentyThree}{95\%}
\providecommand{\acqTalent}{55}
\providecommand{\acqTalentShare}{9\%}
\providecommand{\acqTargets}{200}
\providecommand{\acqTauIres}{24\%}
\providecommand{\acqTenFive}{48\%}
\providecommand{\acqTenMean}{8.0}
\providecommand{\acqTenSqFive}{23\%}
\providecommand{\acqTenTen}{28\%}
\providecommand{\acqTirr}{5}
\providecommand{\acqTlbEur}{EUR 395 million}
\providecommand{\acqTlbUsd}{\$1.25 billion}
\providecommand{\acqTotalAssets}{\$6.98 billion}
\providecommand{\acqTpBe}{13.5\%}
\providecommand{\acqTpBeAnn}{14.4\%}
\providecommand{\acqTpIl}{10.2\%}
\providecommand{\acqTpIm}{10.1\%}
\providecommand{\acqTpImAnn}{10.6\%}
\providecommand{\acqTpIp}{9.1\%}
\providecommand{\acqTpIuSp}{12.4\%}
\providecommand{\acqTpMult}{18.2}
\providecommand{\acqTpPmaxM}{10.6}
\providecommand{\acqTpPmaxP}{4.1}
\providecommand{\acqTpPmaxShare}{54\%}
\providecommand{\acqTpPrice}{18.5}
\providecommand{\acqTpUplift}{106\%}
\providecommand{\acqTpUpliftFast}{101\%}
\providecommand{\acqTractive}{\$759 million}
\providecommand{\acqTractiveWiki}{\$900 million}
\providecommand{\acqTrep}{5}
\providecommand{\acqTrxQOne}{6.3\%}
\providecommand{\acqTrxTwentyFive}{6.5\%}
\providecommand{\acqTrxTwentyFour}{4.8\%}
\providecommand{\acqTtmNi}{\$139.5 million}
\providecommand{\acqTtmRev}{\$1.65 billion}
\providecommand{\acqUcPtsAnnHi}{0.30}
\providecommand{\acqUcPtsHi}{0.59}
\providecommand{\acqUnlForLMax}{31\%}
\providecommand{\acqUnlForLMin}{30\%}
\providecommand{\acqUnpricedDebt}{2{,}234}
\providecommand{\acqVimDays}{327}
\providecommand{\acqVimPre}{\$376.9 million}
\providecommand{\acqVimeoB}{\$1.38 billion}
\providecommand{\acqVimeoPrem}{91\%}
\providecommand{\acqVimeoPx}{\$7.85}
\providecommand{\acqWetPlanned}{75\%}
\providecommand{\acqWetStaff}{350}
\providecommand{\acqXbOne}{2.5}
\providecommand{\acqXbThree}{9}
\providecommand{\acqXbTwo}{5}
\providecommand{\acqXgRatio}{0.51}
% F-1 figures and example values, generated by code/verify_part2_acquirer.py

Sections~\ref{sec:tasks} to \ref{sec:regimes} treat readiness, integration cost, change capacity and learning as attributes of the firm being transformed, and Proposition~\ref{prop:reg-lbo} gives a new owner the target's own capability. A serial acquirer that transforms what it buys breaks both conventions. Its capability moves with it from business to business, so the same process can be worth more to it than to the business itself (Proposition~\ref{prop:acq-owner}); it can replace an organization instead of transforming it in place, escaping Part~I's speed limit at a price (Proposition~\ref{prop:acq-replace}); and holding businesses forever while testing deals on a \acqTirr-year IRR with a terminal value at a cost of capital makes its hurdle a test of the entry price against the value of the business as it will run it (Proposition~\ref{prop:acq-pc}). Retention data constrain the lock-in cap and the customer penalty only in levels (Proposition~\ref{prop:acq-retention}). When the base shrinks and new customers do not offset it, buying growth is a treadmill, and deal size, review of a series of deals and consumer law limit it (Proposition~\ref{prop:acq-limits}). How the score itself should be measured when growth is bought is treated with Part~I (Section~\ref{sec:acquired-growth}). Section~\ref{sec:acq-case} applies the results to Bending Spoons S.p.A.\ \citep{bendingspoons2026f1,bendingspoons2026prospectus,bendingspoons2026sixk}, reporting criticism and the company's claims as each source states them. Figures from its filings (the F-1, the final prospectus and its reports on Form 6-K) are recomputed by \texttt{code/verify\_part2\_acquirer.py} from \texttt{notes/case\_bending\_spoons\_f1\_extract.md}; examples are illustrative and conditional on stated scenario values; proofs are in Appendix~\ref{app:proofs-acquirer}. Annual rates enter as $\ln(1+r)$, and IRRs written $\mathfrak i$ are continuously compounded.

\subsection{Owner-level transformation capability}\label{sec:acq-owner}

\begin{assumption}[Owner capability]\label{ass:acq-owner}
An owner $\mathsf o$ carries a capability $\mathcal U_{\mathsf o}$ to every process of every business it owns: platform terms $(u_F,N_{\rm pl},\mathbb 1^{\rm reb})$, execution $(\varpi^{\mathsf o},T^{\mathsf o}_h)$ and a core team $(N^{\rm core},s^{\rm tech})$. Relative to the transition of Assumption~\ref{ass:tasks-transition} run by the business on its own:
\begin{enumerate}
\item[(i)] (Platform.) The platform already holds the share $u_F\in[0,1]$ of the build and integration cost $F^{\rm bi}:=F-F^{\rm rd}$, where $F^{\rm rd}$ is redesign, legal review, training and change management; the business's $i$-th program in the owner's sequence pays $F^{\mathsf o}:=(1-u_F)F^{\rm bi}i^{-\epsilon_L}+F^{\rm rd}$, with a learning elasticity $\epsilon_L\ge0$ across deals. The run costs of Assumption~\ref{ass:tasks-ancillary}, $C^{\rm run}:=W^0-W$ with $W^0$ the saving before them, consist of maintenance and revalidation $C^{\rm pool}$, the risk cost $C^{\rm risk}(C^{\rm inc})$ of Proposition~\ref{prop:tasks-security}, and the technical debt $C^{\rm spec}$ of integrations into the business's own systems. A platform that serves $N_{\rm pl}$ businesses' versions of the process shares $C^{\rm pool}$ and the risk cost of the pooled exposure among them; $C^{\rm spec}$ is shared only if the business's systems are rebuilt on the platform ($\mathbb 1^{\rm reb}=1$):
\begin{equation}\label{eq:acq-Wo}
W^{\mathsf o}=W^0-\frac{C^{\rm pool}+\mathbb 1^{\rm reb}C^{\rm spec}+C^{\rm risk}(N_{\rm pl}C^{\rm inc})}{N_{\rm pl}}-(1-\mathbb 1^{\rm reb})C^{\rm spec}.
\end{equation}
\item[(ii)] (Execution.) The owner achieves the success probability $\varpi^{\mathsf o}$, for example through the learning of Assumption~\ref{ass:ext-portfolio} with experience transferred across its businesses, and the horizon $T^{\mathsf o}_h$, unbounded for a permanent owner whose platform pays for rebuilds.
\item[(iii)] (Core team.) A team of $N^{\rm core}$ FTE, the share $s^{\rm tech}$ in technical roles, can be assigned without the hiring lead $T^{\rm lead}$ of Assumption~\ref{ass:labor-talent}. Assigning $n^{\rm core}_i$ FTE to business $i$ adds $s^{\rm tech}n^{\rm core}_i$ to its engineering capacity and $(1-s^{\rm tech})n^{\rm core}_i$ to the change managers behind its capacity $K_{\rm chg}$ (Lemma~\ref{lem:ext-staffing}(c)), with $\sum_in^{\rm core}_i\le N^{\rm core}$.
\end{enumerate}
The business on its own, and a sponsor without a platform, have $u_F=0$, $N_{\rm pl}=1$, $\mathbb 1^{\rm reb}=0$, $\varpi^{\mathsf o}=\varpi_p$, $T^{\mathsf o}_h=T_h$ and $N^{\rm core}=0$.
\end{assumption}

\begin{proposition}[The value of owner capability]\label{prop:acq-owner}
Under Assumption~\ref{ass:acq-owner}, let $\mathrm{NPV}_p$ be the business's own value of the transition \eqref{eq:tasks-npv} and $\mathrm{NPV}^{\mathsf o}_p$ its value under $\mathcal U_{\mathsf o}$, with $p_u$ and $\mathcal S_r$ of Proposition~\ref{prop:tasks-transition}, $\mathcal S^{\mathsf o}_r$ at the horizon $T^{\mathsf o}_h$.
\begin{enumerate}
\item[(a)] (Three channels.)
\begin{equation}\label{eq:acq-channels}
\mathrm{NPV}^{\mathsf o}_p-\mathrm{NPV}_p=(F-F^{\mathsf o})+\bar b_p\bigl(\varpi^{\mathsf o}W^{\mathsf o}\mathcal S^{\mathsf o}_r-\varpi_pW\mathcal S_r\bigr)-e^{-rT_{\rm par}}(\varpi^{\mathsf o}-\varpi_p)\bigl[(1-p_u)C^{\rm str}-p_uL_{\rm roll}\bigr].
\end{equation}
\item[(b)] (Scale and rebuilding.) $W^{\mathsf o}$ is nondecreasing in $N_{\rm pl}$ and in $\mathbb 1^{\rm reb}$. Without a rebuild, $W^{\mathsf o}<W^0-C^{\rm spec}$ for every $N_{\rm pl}$, so if $C^{\rm spec}\ge W^0$ no platform scale makes the process save money. With a rebuild, $C^{\rm run}>W^0>0$ and a risk cost that grows less than proportionally to exposure in the limit (as under the exponential residual of Proposition~\ref{prop:tasks-security}(b)), $W^{\mathsf o}>0$ if and only if $N_{\rm pl}>N^{\min}_{\rm pl}$, the unique root above one of $NW^0=C^{\rm pool}+C^{\rm spec}+C^{\rm risk}(NC^{\rm inc})$, and $N^{\min}_{\rm pl}\ge C^{\rm run}/W^0$, strictly if the risk cost is strictly increasing.
\item[(c)] (The transition-cost floor.) At given $\varpi^{\mathsf o}$ and $T^{\mathsf o}_h$, $\mathrm{NPV}^{\mathsf o}_p\le\overline{\mathrm{NPV}}{}^{\mathsf o}_p$, its value at $F^{\mathsf o}=F^{\rm rd}$ and $W^{\mathsf o}=W^0$. If the floor is negative, no platform makes the process pay at that success and horizon: the parallel premium, the productivity dip, stranded costs or rollback must fall instead. $\mathrm{NPV}^{\mathsf o}_p$ is nondecreasing in $u_F$, $N_{\rm pl}$ and $\mathbb 1^{\rm reb}$, in $T^{\mathsf o}_h$ where $W^{\mathsf o}\ge0$, and in $\varpi^{\mathsf o}$ where $\bar b_pW^{\mathsf o}\mathcal S^{\mathsf o}_r\ge e^{-rT_{\rm par}}[(1-p_u)C^{\rm str}-p_uL_{\rm roll}]$; where these hold, the capabilities that make a losing process pay form an upper set.
\item[(d)] (Who values the business most.) In the frictionless benchmark of Proposition~\ref{prop:reg-lbo}(b), with $\Delta V^{\rm cap}_{\mathsf o}:=V^{\rm prog}(\Pi^{\mathsf o};\mathcal U_{\mathsf o})-V^{\rm prog}(\Pi^{\rm pe};\mathcal U_i)$ the value of the owner's plan with its capability over the sponsor's plan with the business's own,
\begin{equation}\label{eq:acq-pmax}
\mathfrak p^{\max}_{\mathsf o}=\mathfrak p^{\max}_{\rm pe}+\Delta V^{\rm cap}_{\mathsf o}/P_0 .
\end{equation}
The business is worth more to the owner than to itself if and only if $V^{\rm prog}(\Pi^{\mathsf o};\mathcal U_{\mathsf o})-V^{\rm prog}(\Pi^{\rm pub};\mathcal U_i)>C_{\rm tx}+e^{-r_wT}p_bL_b$, and more than to a sponsor without a platform if and only if $\Delta V^{\rm cap}_{\mathsf o}>0$. If the business goes to the highest valuation at the second-highest price, the winner keeps only its valuation in excess of the runner-up's: capability that rivals also hold earns its holder nothing.
\item[(e)] (Deal count and the team.) If $u_F$, $N_{\rm pl}$ and $\varpi^{\mathsf o}$ are nondecreasing in the owner's deal count $i$ and the conditions of (c) hold, then $\mathrm{NPV}^{\mathsf o}_p$ of the $i$-th business, and with it $\Delta V^{\rm cap}_{\mathsf o}$ and $\mathfrak p^{\max}_{\mathsf o}$, are nondecreasing in $i$. If every transformation in flight needs at least $\underline n^{\rm core}$ FTE of the team, at most $\lfloor N^{\rm core}/\underline n^{\rm core}\rfloor$ run at once.
\end{enumerate}
\end{proposition}

\noindent\emph{Intuition.} A platform owner owns a cheaper technology of transformation and carries it to each business; auction logic hands the resulting value to the business's holders unless the capability is scarce, like Part~I's rents on scarce inputs (Proposition~\ref{prop:entry}(c)). Pooling helps only with costs that can be pooled: technical debt in the business's own systems survives any scale until the systems are rebuilt. The evidence supports the direction of each channel. Acquirer fixed effects explain about as much of acquirer returns as all firm and deal characteristics together and persist across chief executives \citep{golubov2015extraordinary}; acquirers able to integrate IT across businesses earn higher returns \citep{tanriverdi2011cross}; establishments taken over by U.S.\ multinationals raised the productivity of their IT \citep{bloom2012americans}; acquired dialysis facilities adopted their parent's practices \citep{eliason2020how}; codified acquisition tools, not experience as such, improved bank acquisitions \citep{zollo2004deliberate}; and a high or erratic acquisition rate lowers serial acquirers' performance unless experience offsets it \citep{laamanen2008performance}, the team limit of part (e).

\noindent\emph{Example (illustrative): when would an owner make the invoice chain pay?} The core's running example is worth \acqIcNpv{} on its own, because run costs of \acqIcCrun{} a year exceed the saving of \acqIcWzero{} before them (Section~\ref{sec:tasks-ancillary}). Technical debt in the firm's own integrations, \acqIcCspec, is below the gross saving, so by part (b) pooling alone can make the process save money, but never more than \acqIcLimNoReb{} a year, and only once more than \acqIcNminNoReb{} businesses share its maintenance, revalidation and security; with a rebuild of its systems, it saves money once more than \acqIcNmin{} businesses share it. With all \acqIcFreuse{} of build and integration reused, systems rebuilt and unlimited scale, the floor of part (c) is \acqIcFloorFour{} over the build's life of \acqIcTh{} years, and \acqIcFloorFourCertain{} even with certain success, because parallel running, the productivity dip and stranded and rollback costs remain; with an unlimited horizon it is \acqIcFloorInf{} at the firm's success probability of \acqIcVarpi, which is above the break-even of \acqIcVarpiBe, and \acqIcFloorVone{} with certain success. With certain success and the other levers at those values, the sign flips only above a platform shared by \acqIcNbe{} businesses, a reuse share of \acqIcUFbe{} or a horizon of \acqIcThbe{} years. These are conditions, not a claim about any owner: turning this case around takes capability on several margins at once, or lower transition costs.

\subsection{Replacing instead of transforming in place}\label{sec:acq-replace}

\begin{assumption}[Replacement]\label{ass:acq-replace}
A business runs a program of $N$ identical processes, each saving $W$ a year once embedded (Assumption~\ref{ass:program}). In place, change capacity limits each wave to $n_{\rm cap}$ processes (Proposition~\ref{prop:ext-waves}(d)); waves last $T$, and those of wave $k$ embed at $t_k:=T_{\rm emb}+(k-1)T$, $k=1,\dots,K:=\lceil N/n_{\rm cap}\rceil$. The owner can instead replace: at date 0 it starts the procedures of Assumption~\ref{ass:labor-proc} for the share $\mathsf d^{\rm rpl}$ of the legacy staff in each jurisdiction $\mathfrak c$, its core team rebuilds all $N$ processes on its platform in $T^{\rm bld}$, and all cut over at $t^{\rm rpl}:=\max\{T^{\rm bld},\max_{\mathfrak c}T^{\rm leg}_{\mathfrak c}\}$, after which the legacy organization is gone. Relative to transforming in place, replacing:
\begin{enumerate}
\item[(i)] adds severance and procedure costs $\Delta C^{\rm sev}$ (Assumption~\ref{ass:labor-proc}(ii), with redeployment $u^{\rm rpl}_d$ near zero);
\item[(ii)] lowers the context stock of Assumption~\ref{ass:ext-shocks}(iii) by $(1-\kappa_{\rm doc})\mathsf d^{\rm rpl}$ at $t^{\rm rpl}$, recovering at $\eta_K$, at the cost $C^{\rm ctx}:=L^{\rm ctx}(1-\kappa_{\rm doc})\mathsf d^{\rm rpl}e^{-rt^{\rm rpl}}/(r+\eta_K)$;
\item[(iii)] lowers retention growth by $\delta^G$ for $T^{\rm cu}$ (Assumption~\ref{ass:ext-theta}(v)), a permanent loss of the share $1-e^{-\delta^GT^{\rm cu}}$ of revenue, at the cost $C^{\rm cu}:=X^{\mathsf o}S(1-e^{-\delta^GT^{\rm cu}})e^{-rt^{\rm rpl}}$ at the owner's growth weight $X^{\mathsf o}$ (Proposition~\ref{prop:dynamic});
\item[(iv)] cuts over its $n_{\rm fin}$ financial-reporting processes together at $t^{\rm rpl}$, whose position in the fiscal year is $z^{\rm rpl}$, while in place each wave cuts over right after a close. A process cut over with $H$ to go to year-end, among $n$ handled at once, leaves an unremediated material weakness at that year-end with probability $P_{\rm mw}(H;n)$, equation \eqref{eq:acct-pmw} evaluated at $p_{\rm def}(n)$, $T_{\rm dc}(n)$ and $T_{\rm rem}(n)$, each nondecreasing in $n$ when control owners congest. Deficiencies are independent across processes. A year with at least one weakness costs $L_{\rm mw}$, since management may not conclude that ICFR is effective if one or more material weaknesses exist (Item 308(a)(3), \citealp{cfr2026sec229308}); the expected excess over the in-place schedule is $\Delta C^{\rm mw}$; and replacing the business's control layers by the owner's design changes its floor \eqref{eq:acct-amin} by $\Delta a_{\min}$;
\item[(v)] loses $L^{\rm rpl}_{\rm fail}$ on failure, probability $1-\varpi^{\rm rpl}$, since no rollback is possible;
\item[(vi)] changes one-time costs by $\Delta F$ (the core team's rebuild less the in-place builds) and avoids the dual-running cost $\Delta C^{\rm par}$ of the in-place waves: until dismissal, legacy staff keep the old processes running, the free parallel running of Proposition~\ref{prop:labor-lead}(e).
\end{enumerate}
\end{assumption}

\begin{proposition}[Replace or transform in place]\label{prop:acq-replace}
Under Assumption~\ref{ass:acq-replace}:
\begin{enumerate}
\item[(a)] (Condition.) Replacing beats transforming in place if and only if $\mathcal G^{\rm spd}>\Delta\mathcal C$, where
\begin{equation}\label{eq:acq-replace}
\begin{gathered}
\mathcal G^{\rm spd}:=\frac Wr\sum_{k=1}^Kn_k\bigl(e^{-rt^{\rm rpl}}-e^{-rt_k}\bigr),\\
\Delta\mathcal C:=\Delta C^{\rm sev}+C^{\rm ctx}+C^{\rm cu}+\Delta C^{\rm mw}+\frac{S\,\Delta a_{\min}}r+(1-\varpi^{\rm rpl})e^{-rt^{\rm rpl}}L^{\rm rpl}_{\rm fail}+\Delta F-\Delta C^{\rm par}.
\end{gathered}
\end{equation}
\item[(b)] (What favors replacing.) $\mathcal G^{\rm spd}$ falls with $t^{\rm rpl}$, hence with the longest legal lead, and with $n_{\rm cap}$. $\Delta\mathcal C$ rises with severance, with $\delta^G$, $T^{\rm cu}$ and $X^{\mathsf o}$, and with $L^{\rm ctx}$, and falls with $\kappa_{\rm doc}$. Tight in-place capacity, short leads, low severance, documented context, sticky customers and a low growth weight favor replacing.
\item[(c)] (Speed gain.) With full waves, $\mathcal G^{\rm spd}=(n_{\rm cap}W/r)\bigl[Ke^{-rt^{\rm rpl}}-e^{-rT_{\rm emb}}(1-e^{-rKT})/(1-e^{-rT})\bigr]$, and as $r\to0$, $\mathcal G^{\rm spd}\to NW(\bar t^{\rm ip}-t^{\rm rpl})$, with $\bar t^{\rm ip}$ the mean in-place embedding date.
\item[(d)] (Pace and material weaknesses.) At equal positions in the fiscal year and without congestion, bunching $n$ cutovers into one year does not raise the expected number of years with an unremediated weakness, since $1-(1-P)^n\le nP$. If $nP_{\rm mw}(1-T_{\rm cl};1)<1$, bunching at $t^{\rm rpl}$ costs more than spreading the cutovers right after $n$ closes if and only if
\begin{equation}\label{eq:acq-mw}
P_{\rm mw}(1-z^{\rm rpl};n)>1-\bigl[1-nP_{\rm mw}(1-T_{\rm cl};1)\bigr]^{1/n},
\end{equation}
which congestion (through $p_{\rm def}$, $T_{\rm dc}$ and $T_{\rm rem}$) or a late cutover date (a short $1-z^{\rm rpl}$) can produce. If each weakness costs $L_{\rm mw}$ instead (Corollary~\ref{cor:acct-npv}), bunching and spreading cost the same without congestion, and any congestion makes bunching costlier.
\item[(e)] (Attainable frontier.) Let each process weakly expand the feasible set, and let $\delta\theta^{\rm rpl}$ be the permanent parameter change replacement causes (a lasting fall in $G_i$, or in $m_0$ through lost context). Replacement's attainable set is $\mathcal F(\theta_0)$ before $t^{\rm rpl}$ and $\mathcal F(\theta_{\mathcal P}+\delta\theta^{\rm rpl})$ after, free of the speed limit of Proposition~\ref{prop:attainable}(c). If $\delta\theta^{\rm rpl}=0$, it contains the in-place set $\mathcal A_t$ at every $t\ge t^{\rm rpl}$. If $\delta\theta^{\rm rpl}$ lowers $R^*_X$, replacement attains the higher score from $t^{\rm rpl}$ until the first date at which the processes embedded in place raise $R^{\rm att}_X$ above $R^*_X(\theta_{\mathcal P}+\delta\theta^{\rm rpl})$, and the lower one afterward. All processes then run in one wave on one platform, so the probability that none embeds is at least $1-p_P$ (Proposition~\ref{prop:ext-portfolio}(c)).
\end{enumerate}
\end{proposition}

\noindent\emph{Intuition.} Part~I's pace results limit how fast an organization can change itself. An owner that dismisses most of the legacy organization and rebuilds with its own team need not absorb the change, so the limit becomes the build time and the legal leads; the price is paid once, in severance, context, customers, control risk and the lack of a fallback. A single cut also fits Proposition~\ref{prop:labor-survivor}(b). Acquirers do restructure fast: full-firm acquirers sell 27\% and close 19\% of target plants within three years, and retained plants gain productivity \citep{maksimovic2011post}. Labor law prices this path: major increases in employment protection reduce takeovers by 14 to 27\% and their synergies by more than half \citep{dessaint2017employment}. Downsizing harms reputation \citep{love2009character}, displaced target employees bear lasting earnings losses \citep{lagaras2026mas}, and replacing an acquired bank's top managers lowered performance \citep{zollo2004deliberate}. Material weaknesses are more common in firms growing rapidly or restructuring \citep{doyle2007determinants} and after organizational changes \citep{ashbaughskaife2007discovery}; part (d) shows that bunching at a given date does not raise the risk, so integration risk needs congested control owners, lost context, or a cutover date that the build and the legal leads place late in the fiscal year.

\noindent\emph{Example (illustrative).} In the twelve-process program of Proposition~\ref{prop:ext-waves} (\$1 million a year per process, six-month waves, four a wave, embedding at 0.5, 1.0 and 1.5 years), at $r_w=\acqRtv$, a three-month rebuild and the labor section's legal leads for 150 redundancies (Proposition~\ref{prop:labor-lead}), replacing brings savings forward by \$\acqRpGainUS{} million in the United States, the United Kingdom and India, where the build binds, \$\acqRpGainDE{} million in Germany and \$\acqRpGainFR{} million in France (\$\acqRpGainHalf{} million everywhere with a six-month rebuild). That gain is the budget for $\Delta\mathcal C$: at the parameter file's central severance and a \$100{,}000 loaded wage it pays for about \acqRpFteUS{} (U.S.), \acqRpFteDE{} (Germany), \acqRpFteFR{} (France) or \acqRpFteUK{} (U.K.) extra displaced FTE, before other costs; on \$500 million of revenue it covers a permanent revenue loss of $\mathcal G^{\rm spd}/(X^{\mathsf o}S)$: \acqRpCustOne{} at a growth weight of one, and \acqRpCustXbOne, \acqRpCustXbTwo{} or \acqRpCustXbThree{} at illustrative weights of \acqXbOne, \acqXbTwo{} or \acqXbThree. % XBOUNDS: placeholder weights (macros acqXbOne to acqXbThree, set in code/verify_part2_acquirer.py)
The break-even falls as $1/X^{\mathsf o}$ and never changes sign, so at any weight speed alone pays only for small permanent losses of revenue. With the accounting section's illustrative values (Proposition~\ref{prop:acct-cutover}), four cutovers right after a close give a \acqMwBunch{} risk of a weakness in the year against \acqMwSpread{} expected weakness-years if spread right after four closes, and congestion that lifts the per-cutover risk by more than \acqMwCrossRise{} of its value, from \acqMwOneExact{} to \acqMwCross, reverses the ranking. So does the date alone: each cutover carries \acqMwOneExact{} right after a close but \acqMwLate{} with half a year or less to go, four bunched there give \acqMwBunchLate, and bunching at one date is worse than spreading once that date falls later than \acqMwZ{} of the fiscal year, about \acqMwZdays{} days after the earliest open date. Processes cut over together onto the owner's platform share its ITGCs (Proposition~\ref{prop:acct-itgc}(b)), so their deficiencies are correlated; that leaves the direction of part (d) unchanged under the cost of a year with a weakness but makes a weakness pervasive. If management excludes an acquired business from its first assessment, as SEC staff allow for up to a year \citep[Question~3]{sec2007faq}, a replacement cut over within that year is first assessed at a later year-end, so $P_{\rm mw}$ falls: the exclusion favors cutting over early in the first year after an acquisition. The company's reorganization-related expense, which its non-GAAP measures add back, was \acqReorgTwentyFive{} of revenue in 2025 and \acqReorgQOne{} in the first quarter of 2026 \citep{bendingspoons2026prospectus}, and \acqQTwoReorgAbs{} (\acqQTwoReorgX) in the second, reflecting separation packages at AOL, Eventbrite, Tractive and Vimeo \citep{bendingspoons2026sixk}. Because the businesses reorganized are part of the total, the cost per dollar of their revenue is higher still. Scaled to the example's \$500 million of revenue, 2025's share is at least \acqReorgUsdEx, \acqReorgRatio{} times the U.S.\ speed gain: at that cost speed alone does not justify replacing, which must also deliver a cheaper rebuild ($\Delta F<0$) or a permanent change in its favor in part (e).

\subsection{A permanent-capital serial acquirer}\label{sec:acq-pc}

\begin{assumption}[Permanent capital]\label{ass:acq-pc}
The regime $\mathsf o=\mathrm{pc}$ joins those of Assumption~\ref{ass:reg-regime}:
\begin{enumerate}
\item[(i)] (Horizon.) $T_{\rm pc}=\infty$: businesses are never sold, and a plan's value is \eqref{eq:reg-value} with no terminal term.
\item[(ii)] (Hurdle test.) A business bought at enterprise value $P$ (costs included) is acquired only if its method IRRs clear $\bar{\mathfrak i}_{\rm U}$ on unlevered and $\bar{\mathfrak i}_{\rm L}$ on levered cash flows, where the method IRR $\mathfrak i^{\rm m}$ solves
\begin{equation}\label{eq:acq-method}
P=\int_0^{T^{\rm IRR}}e^{-\mathfrak it}\mathrm{FCF}(t)\,dt+e^{-\mathfrak iT^{\rm IRR}}\mathrm{TV},
\end{equation}
$\mathrm{FCF}$ is the business's free cash flow under the owner's plan, and $\mathrm{TV}$ values later flows at the rate $r_{\rm tv}$ (recent deals) or at an EBITDA multiple (earlier deals). The permanent IRR $\mathfrak i^\infty$ solves $P=\int_0^\infty e^{-\mathfrak it}\mathrm{FCF}(t)\,dt$.
\item[(iii)] (Debt.) Acquisition debt is $\mathrm{ND}_0=\min\{\overline{\mathrm{LTV}}\,P,\mathrm{ND}^{\rm rep}\}$, where $\mathrm{ND}^{\rm rep}:=\int_0^{T^{\rm rep}}e^{-r_{\rm debt}t}\mathrm{FCF}(t)\,dt$ is the debt that projected free cash flow repays within $T^{\rm rep}$; free cash flow sweeps the debt until it is repaid; and a consolidated covenant $\mathrm{ND}_t\le\overline{\rm lev}^{\rm pc}\,\mathrm{EB}_t$ applies as in Proposition~\ref{prop:reg-dip}(a).
\item[(iv)] (Pressure.) The myopia weight is $\mathfrak w_{\rm pc}\in[0,\mathfrak w_{\rm dc}]$: private before listing, then public with dual-class shares (Proposition~\ref{prop:reg-hybrid}(a)).
\item[(v)] (Capability.) $\mathcal U_{\rm pc}$ of Assumption~\ref{ass:acq-owner}.
\end{enumerate}
The F-1 sets $T^{\rm IRR}=T^{\rm rep}=\acqTirr$ years, $\bar{\mathfrak i}_{\rm U}=\ln(1+\acqHurU)=\acqIrrU$, $\bar{\mathfrak i}_{\rm L}=\ln(1+\acqHurL)=\acqIrrL$, $\overline{\mathrm{LTV}}=\acqLtv$ and $\overline{\rm lev}^{\rm pc}=\acqCov$, against an actual \acqLevTwentyFive{} at end-2025 and \acqLevQOne{} at end-Q1 2026 (grounding a); the \acqLevQTwo{} reported for mid-2026 uses an EBITDA that includes achieved and certain expected cost savings from reorganizations \citep{bendingspoons2026sixk}, a looser definition than the model's $\mathrm{EB}_t$. We set $r_{\rm debt}=\ln(1+\acqNotesUsdC)=\acqRdebt$, the coupon on the dollar tranche of the senior secured notes that two of its subsidiaries priced on 28 October 2026 (\acqNotesEurC{} on the euro tranche; a coupon, not a yield; the company does not issue or guarantee them), as its marginal cost of debt (a). The filings give no $r_{\rm tv}$; the examples use the software cost of capital $r_w=\acqRtv$ (b).
\end{assumption}

\begin{proposition}[Permanent capital, the hurdle and leverage]\label{prop:acq-pc}
Under Assumption~\ref{ass:acq-pc}, let $\mathrm{FCF}\ge0$ after date 0, let $\mathrm{TV}$ be valued at $r_{\rm tv}$, and let $V^{\rm tv}:=\int_0^\infty e^{-r_{\rm tv}t}\mathrm{FCF}(t)\,dt$ be the value of the plan at $r_{\rm tv}$.
\begin{enumerate}
\item[(a)] (A $T^{\rm IRR}$-year hold sold at $r_{\rm tv}$.) $\mathfrak i^{\rm m}\ge\mathfrak i^\infty$ if and only if $\mathfrak i^\infty\ge r_{\rm tv}$. For the hurdle, the permanent owner is a holder for $T^{\rm IRR}$ years who exits at a valuation at $r_{\rm tv}$: the market rule of Assumption~\ref{ass:reg-regime}(ii) with the buyer's rate $r_{\rm tv}$. The never-sell policy changes realized flows, not the test.
\item[(b)] (Entry price.) $\mathfrak i^{\rm m}\ge\bar{\mathfrak i}$ if and only if $P\le P^{\max}(\bar{\mathfrak i})$, the right side of \eqref{eq:acq-method} at $\bar{\mathfrak i}$, and for $\bar{\mathfrak i}\ge r_{\rm tv}$
\begin{equation}\label{eq:acq-pmaxbound}
e^{-(\bar{\mathfrak i}-r_{\rm tv})T^{\rm IRR}}\,V^{\rm tv}\le P^{\max}(\bar{\mathfrak i})\le V^{\rm tv}.
\end{equation}
If $P=(1+\mathfrak p)V^{\rm tv}_{\rm sq}$ with $V^{\rm tv}=V^{\rm tv}_{\rm sq}+V^{\rm tv}_{\rm prog}$, the hurdle holds if $V^{\rm tv}_{\rm prog}/V^{\rm tv}_{\rm sq}\ge(1+\mathfrak p)e^{(\bar{\mathfrak i}-r_{\rm tv})T^{\rm IRR}}-1$, and only if $V^{\rm tv}_{\rm prog}/V^{\rm tv}_{\rm sq}\ge\mathfrak p$.
\item[(c)] (Leverage.) Let $0<\mathrm{ND}_0<P$. (i) The levered method IRR exceeds the unlevered one if and only if the unlevered one exceeds $r_{\rm debt}$. (ii) If $\mathrm{ND}^{\rm rep}\le\overline{\mathrm{LTV}}P$ and $T^{\rm rep}=T^{\rm IRR}$, equity receives nothing before $T^{\rm IRR}$ and $\mathfrak i^{\rm m}_{\rm L}=\ln[\mathrm{TV}/(P-\mathrm{ND}^{\rm rep})]/T^{\rm IRR}$.
\item[(d)] (Growth weight.) If the owner holds internal decisions to $\bar{\mathfrak i}_{\rm U}$, its weight on a permanent dollar of revenue with free-cash-flow margin $\hat m$ that grows at the retention rate $G_i$ is $X^{G}=\hat m/(\bar{\mathfrak i}_{\rm U}-G_i)$ as a going concern (Proposition~\ref{prop:reg-weight}(c) as the exit recedes), $(r_w-G_i)/(\bar{\mathfrak i}_{\rm U}-G_i)$ times that of an owner at $r_w$.
\end{enumerate}
\end{proposition}

\noindent\emph{Intuition.} Valuing flows after $T^{\rm IRR}$ at a cost of capital below the hurdle credits the owner with a sale at that valuation, so the hurdle tests the entry price against the value of the business as the owner will run it. A lower price relative to status-quo value, a program worth a multiple of that value (capability) and earlier cash flow inside the window (speed) all raise the admissible price; speed moves it from the lower toward the upper end of \eqref{eq:acq-pmaxbound}. Leverage helps only deals that already beat the debt rate.

\noindent\emph{Example (illustrative): the take-private firm of Proposition~\ref{prop:reg-lbo}.} At the sponsor's price, \acqTpPrice{} times EBITDA with costs, the permanent-capital owner running the same program earns an unlevered method IRR of \acqTpIm{} (\acqTpImAnn{} annual) and \acqTpIl{} levered under its debt policy, below the sponsor's levered break-even of \acqTpBe{} and the sponsor deal's unlevered \acqTpIuSp{} (which exits at the multiple paid, above the value at $r_w$); its unlevered hurdle of \acqIrrU{} exceeds the sponsor's levered survey target of \acqIrrPe. Clearing it needs a price of at most \acqTpPmaxM{} times EBITDA (\acqTpPmaxShare{} of the plan's value at $r_w$; lower bound \acqBoundFac), or a run-rate uplift of \acqTpUplift{} of entry EBITDA instead of 30\% (\acqTpUpliftFast{} if reached in six months). The hurdle is met by buying businesses whose status-quo value is low relative to what the owner can make of them, not by doing a sponsor's deal faster. A business with no free cash flow in its first year, a margin $f_1$ afterward and revenue decaying at 2025's portfolio retention (\acqNrrTwentyFive) clears it at up to \acqGridMthreeA, \acqGridMfourA{} or \acqGridMfiveA{} times revenue for $f_1$ of 0.3, 0.4 or 0.5 (\acqGridMfourB{} for 0.4 at 2024's \acqNrrTwentyFour); Brightcove's reported price was \acqBcPs{} times its 2024 revenue guidance \citep{streamingmedia2024brightcove}. At these prices the levered method IRR is \acqGridLevMin{} to \acqGridLevMax, below $\bar{\mathfrak i}_{\rm L}$, and for flat-margin businesses the levered hurdle needs unlevered IRRs of \acqUnlForLMin{} to \acqUnlForLMax: if both hurdles are minimums, the levered one binds at the notes' coupon. The final prospectus prices the three deals in its pro forma statements at \acqKEb{} (Eventbrite), \acqKAol{} (AOL) and \acqKVim{} (Vimeo) times annual revenue \citep{bendingspoons2026prospectus}: Eventbrite lies inside the grid's admissible range, which reaches \acqGridMax{} times, and AOL and Vimeo lie above it, so they clear the hurdle only with higher margins sooner (AOL's own 2025 operating margin was \acqAolOpm), better retention, or a lower hurdle. Part (d), for subscription revenue retained at 2025's portfolio NRR, puts the owner's weight on revenue at \acqXgRatio{} of a going concern's at $r_w$.

\subsection{Repricing and retention}\label{sec:acq-retention}

\begin{proposition}[Lock-in, repricing and retention]\label{prop:acq-retention}
\begin{enumerate}
\item[(a)] (Heterogeneous lock-in.) A customer with bill $B$ and a switching cost $C^{\rm sw}$ per contract term $T_c$ stays after a price rise $\upsilon_{\rm rp}$ if and only if $\upsilon_{\rm rp}\le C^{\rm sw}/(T_cB)$, by the cap of Proposition~\ref{prop:tasks-buildbuy}(c). If $\mathsf S_{\rm sw}(x)$ is the share of revenue whose ratio $C^{\rm sw}/(T_cB)$ exceeds $x$, a uniform repricing retains revenue $(1+\upsilon_{\rm rp})\mathsf S_{\rm sw}(\upsilon_{\rm rp})$. With one ratio for all customers, retention after a repricing is either $1+\upsilon_{\rm rp}$ or zero, so retention strictly between zero and one after a large repricing requires heterogeneous switching costs. The revenue-maximizing repricing solves the inverse-hazard condition
\begin{equation}\label{eq:acq-reprice}
\mathsf S_{\rm sw}(\upsilon)=-(1+\upsilon)\,\mathsf S'_{\rm sw}(\upsilon),
\end{equation}
and with exponential ratios of mean $\bar x>1$ it is $\bar x-1$.
\item[(b)] (What NRR identifies.) $\mathrm{NRR}=(1+\bar\upsilon)(1-\bar c)(1+\bar e)$ for the average price change, revenue churn and expansion. The customer penalty of Assumption~\ref{ass:ext-theta}(v) is the change $\delta^G$ in $G_i$, which for a business whose revenue is all subscription revenue is $\ln\mathrm{NRR}^{\rm pre}-\ln\mathrm{NRR}^{\rm post}$, and it needs retention before the owner acted; levels identify only the zero-spend growth of subscription revenue, $\ln\mathrm{NRR}$, which is $G_i$ of Proposition~\ref{prop:acquired-score}(d) for such a business.
\item[(c)] (Tenure.) If new-customer revenue enters at a rate growing at $g^{\rm in}$ and each cohort's revenue decays at the constant rate $\nu^{\rm rv}$, the share of revenue from customers of tenure at least $\tau$ is $\mathrm{Sh}(\tau)=e^{-(\nu^{\rm rv}+g^{\rm in})\tau}$ and the revenue-weighted mean tenure is $1/(\nu^{\rm rv}+g^{\rm in})$. $\mathrm{Sh}(2\tau)>\mathrm{Sh}(\tau)^2$ is inconsistent with any constant rate: the revenue hazard falls with tenure, or the inflow varied over time.
\item[(d)] (The owner's weight.) For an asset-light business ($k=0$) the customer penalty $\mathcal C$ of \eqref{eq:ext-wedges} is proportional to the growth weight; an owner with the weight of Proposition~\ref{prop:acq-pc}(d) charges a given retention loss $(r_w-G_i)/(\bar{\mathfrak i}_{\rm U}-G_i)$ times what a going concern at $r_w$ with the same margin would.
\end{enumerate}
\end{proposition}

\noindent\emph{Calibration against the case.} Business-level NRR puts the zero-spend growth of each business's subscription revenue, $\ln\mathrm{NRR}$, at \acqGiAol{} (AOL), \acqGiEvernote{} (Evernote), \acqGiRemini{} (Remini) and \acqGiStreamyard{} (StreamYard): every subscription base shrinks without new customers. These are $G_i$ of Proposition~\ref{prop:acquired-score}(d) only for a business whose revenue is all subscription revenue. These are levels; with no pre-acquisition retention in the extract, $\delta^G$ is not identified. In Q1 2026, \acqTenFive{} of subscription revenue came from customers of at least five years and \acqTenTen{} from at least ten, with a revenue-weighted mean tenure of \acqTenMean{} years. Through part (c) these imply $\nu^{\rm rv}+g^{\rm in}$ of \acqHazFive, \acqHazTen{} and \acqHazMean; $\mathrm{Sh}(10)=\acqTenTen$ exceeds $\mathrm{Sh}(5)^2=\acqTenSqFive$, so the hazard falls with tenure, as lock-in that deepens with tenure would produce. NRR implies cohort decay of \acqNuTwentyFive{} (2025) to \acqNuTwentyFour{} (2024), leaving inflow growth of about \acqGinLo{} to \acqGinHi; a stationary two-type fit (\acqMixW{} of revenue at a hazard of \acqMixFast, the rest at \acqMixSlow) matches the tenure moments but implies one-year retention of \acqMixNrr, below every reported NRR, so the data fit only inflows that changed over time. Under the cap of part (a), the customer in the reported Harvest repricing, from about \acqHarvestOld{} to \acqHarvestNew{} a month \citep{subscriptioninsider2026harvest}, stays only if its switching cost per term exceeds about \acqHarvestMult{} times its old bill; whether it stayed is not reported. By part (d), the permanent-capital owner weights a retention loss at \acqXgRatio{} of a going concern's at $r_w$, so it rationally accepts losses that owner would not. Limits: survivorship, since the businesses were bought for durable bases and long-tenure customers are survivors \citep{brown1992survivorship}; aggregation, since portfolio NRR can rise as weights shift with acquisitions; and timing, since the Q1 2023 to Q1 2026 averages include pre-ownership quarters for later acquisitions such as AOL (closed 2 January 2026), unless the F-1 defines them otherwise.

\subsection{Limits to buying growth}\label{sec:acq-limits}

\begin{assumption}[Limits]\label{ass:acq-limits}
The owner of Assumption~\ref{ass:acq-pc} targets revenue growth $g$, of which $g^{\rm org}$ is organic and $g^{\rm acq}=g-g^{\rm org}$ is acquired at $k^{\rm acq}$ per dollar of revenue (Definition~\ref{def:acquired-growth}).
\begin{enumerate}
\item[(i)] (Funding.) Acquisitions are paid from the free-cash-flow margin $m$ before interest and from net debt held at $\mathrm{nd}<k^{\rm acq}$ per dollar of revenue at the rate $r_{\rm debt}$, with no new equity.
\item[(ii)] (Deal flow.) At most $n^{\rm dl}$ deals close a year (the team limit of Proposition~\ref{prop:acq-owner}(e) over a transformation's duration). A target with revenue $S^{\rm tg}$ offers hurdle headroom $\mathsf h(S^{\rm tg}):=P^{\max}/P=\mathsf h_0(S^{\rm tg}/S^{\rm tg}_0)^{-\epsilon_{\rm sz}}$, $\epsilon_{\rm sz}\ge0$, and clears the hurdle of Proposition~\ref{prop:acq-pc}(b) if and only if $\mathsf h\ge1$.
\item[(iii)] (Review.) Each deal in niche $j$ adds $\Delta s$ to the owner's share there; the $n$-th has net value $V_n-P_n$, nonincreasing in $n$, and is challenged with probability $p^{\rm ch}(n\Delta s)$ if the agency may examine the series and $p^{\rm ch}(\Delta s)$ if it judges each deal alone, with $p^{\rm ch}$ nondecreasing, at a loss $L^{\rm ch}$. A foreign-investment review delays closing by $T^{\rm fdi}$ with probability $p^{\rm fdi}$.
\item[(iv)] (Consumer law.) Consumer or subscription law caps a repricing at $\upsilon^{\rm cp}$, or scales every switching cost by $h^{\rm cp}\in(0,1]$.
\end{enumerate}
\end{assumption}

\begin{proposition}[Treadmill, deal size, review and consumer law]\label{prop:acq-limits}
Under Assumption~\ref{ass:acq-limits}:
\begin{enumerate}
\item[(a)] (Treadmill.) The owner finances growth $g$ if and only if $g\le g^{\max}:=(m-r_{\rm debt}\mathrm{nd}+k^{\rm acq}g^{\rm org})/(k^{\rm acq}-\mathrm{nd})$, which is $g^{\rm org}+m/k^{\rm acq}$ without debt and falls with $k^{\rm acq}$ if and only if $m-r_{\rm debt}\mathrm{nd}+\mathrm{nd}\,g^{\rm org}>0$. Holding growth $g$ requires acquired growth $g+d^{\rm org}$, the target growth plus the organic decay $d^{\rm org}:=-g^{\rm org}$, at a cash cost of $k^{\rm acq}(g+d^{\rm org})$ of revenue a year. When no new customers arrive, $d^{\rm org}=-G_i$ of Proposition~\ref{prop:acquired-score}(d), which is $-\ln\mathrm{NRR}$ when all revenue is subscription revenue; standing still ($g=0$) then requires acquired growth $-G_i$.
\item[(b)] (Deal size.) Holding $g$ requires a mean target revenue of at least $(g-g^{\rm org})S/n^{\rm dl}$, proportional to $S$. If $\epsilon_{\rm sz}>0$, targets clear the hurdle only up to $S^{\rm tg}_{\max}:=S^{\rm tg}_0\mathsf h_0^{1/\epsilon_{\rm sz}}$, so $g$ can be held only while $S\le S^{\max}:=n^{\rm dl}S^{\rm tg}_{\max}/(g-g^{\rm org})$; beyond it, growth is at most $g^{\rm org}+n^{\rm dl}S^{\rm tg}_{\max}/S$, falling toward $g^{\rm org}$, unless the hurdle, the team or the deal rate changes. If $\epsilon_{\rm sz}=0$ and $\mathsf h_0\ge1$, only funding and the supply of targets bind.
\item[(c)] (Series review.) The owner makes the $n$-th deal in a niche if and only if $V_n-P_n\ge p^{\rm ch}(n\Delta s)L^{\rm ch}$ under series review, and stops no later than under deal-by-deal review, so it accumulates a smaller share of the niche; with $p^{\rm ch}$ convex, each further deal's expected regulatory cost rises faster. A foreign-investment review scales a deal's net value by $1-p^{\rm fdi}(1-e^{-rT^{\rm fdi}})$.
\item[(d)] (Consumer law.) A repricing $\upsilon\le\upsilon^{\rm cp}$ retains revenue $(1+\upsilon)\mathsf S_{\rm sw}(\upsilon/h^{\rm cp})$, and the revenue-maximizing repricing of Proposition~\ref{prop:acq-retention}(a) maximizes it; with exponential ratios of mean $\bar x$ it is $\min\{(h^{\rm cp}\bar x-1)^+,\upsilon^{\rm cp}\}$. Retained revenue at every repricing falls as $h^{\rm cp}$ falls: the law tightens the lock-in cap.
\end{enumerate}
\end{proposition}

\noindent\emph{Grounding.} Rates here are continuously compounded. With no new customers and all revenue on subscription, the portfolio's NRR range (\acqNrrLo{} to \acqNrrHi) puts the organic decay at \acqStandHi{} to \acqStandLo{} a year (\acqDeclLo{} to \acqDeclHi{} as annual rates). Standing still then means acquiring \acqStandUsdHi{} to \acqStandUsdLo{} of run-rate a year at 2025 revenue, a cash cost of \acqStandCashLo{} to \acqStandCashHi{} of revenue at the prices of the three deals in the final prospectus, \acqKLo{} to \acqKHi{} times revenue; holding 2025's growth of \acqLogGrowthTwentyFive{} means acquiring \acqReqLo{} to \acqReqHi{} of revenue a year, at a cash cost of \acqReqCashLo{} to \acqReqCashHi{} of revenue at those prices. Revenue that is not subscription revenue and does not decay as fast lowers the decay. New customers offset it as well: in the second quarter of 2026 the company put its organic growth at \acqQTwoOrg{} (\acqOrgLogQTwo{} continuously compounded) on its own unaudited measure \citep{bendingspoons2026sixk}, at which holding 2025's growth needs \acqReqOrg{} a year. Without debt, with $m$ of \acqMgridLo{} to \acqMgridHi{} (scenario values) and $k^{\rm acq}$ at those prices, $g^{\max}$ is \acqGmaxLo{} to \acqGmaxHi{} at 2025's NRR with no new customers and \acqGmaxOrgLo{} to \acqGmaxOrgHi{} at the company's organic growth, far below 2025's growth of \acqLogGrowthTwentyFive; financing inflows of \acqCffTwentyFive{} (\acqCffShare{} of revenue) made up the difference. The size elasticity $\epsilon_{\rm sz}$ is a scenario parameter that the literature signs but does not size. Frequent bidders gain on private and subsidiary targets and lose on public ones \citep{fuller2002what}; acquirers underperform especially in public takeovers \citep{renneboog2019failure}; returns fall from deal to deal \citep{billett2008overconfident}, which rational learning predicts as well as hubris \citep{aktas2009learning}; the persistence of U.S.\ serial acquirers' returns is absent elsewhere \citep{karolyi2025serial}; and add-ons are more frequent and productive under experienced sponsors and larger platforms \citep{hammer2017inorganic}. A practitioner study finds that firms making many small deals earned higher excess returns than peers in 2007--2017 \citep{mckinsey2019lots}. The case moved toward larger targets: Vimeo (\acqPxVim), Brightcove (\acqBcPrice, reported) and Eventbrite (\acqPxEb) were take-privates, AOL (\acqPxAol) was bought from Yahoo \citep{bendingspoons2026prospectus,ffxnow2026aol}, and after the listing it bought Tractive (enterprise value \acqTractive) and the private Airtable (\acqAirEv) and Miro (\acqMiroEv), both completed \citep{bendingspoons2026sixk}. The notes priced in October 2026 were its first bonds, against the F-1's statement that it had issued none. Whether its hurdles fall with scale, as the F-1 is reported to say, needs the filing's wording, which our extract lacks.

\noindent\emph{Regulation.} In the United States the 2023 Merger Guidelines remain in effect \citep{ferguson2025memo}; Guideline 8 is headed ``When a Merger Is Part of a Series of Multiple Acquisitions, the Agencies May Examine the Whole Series'' \citep{doj2023guidelines}, the series review of part (c). The agencies asked for public comment on serial acquisitions and roll-ups in May 2024 \citep{thompsonhine2024rfi}. In the case against U.S.\ Anesthesia Partners, the claims against its private-equity founder, Welsh Carson, were dismissed, and Welsh Carson later accepted a consent order requiring prior approval of anesthesiology acquisitions \citep{stblaw2025welshcarson}; a confidential settlement with the company was reported in April 2026 \citep{insurancejournal2026usap}. The expanded premerger form of 2024 was vacated in February 2026 \citep{bakerbotts2026hsr}. Size thresholds matter for serial buyers: when the U.S.\ threshold rose, investigations of newly exempt deals fell to nearly zero and mergers between competitors rose \citep{wollmann2019stealth}, and about 40\% more deals than expected fall just below the thresholds, followed by higher margins and prices \citep{kepler2023stealth}. In the European Union, the Court of Justice held in \emph{Illumina/Grail} on 3 September 2024 that a member state without jurisdiction over a deal cannot refer it to the Commission \citep{gtlaw2024illumina}, but nine member states can call in deals below their thresholds \citep{whitecase2025callin}; a law-firm summary of the Commission's draft merger guidelines of 30 April 2026 does not mention serial acquisitions or roll-ups \citep{macfarlanes2026draft}. Repricing after a merger can draw competition scrutiny too: five European trade groups, CISPE among them, asked the Commission on 10 July 2026 for interim measures against Broadcom, which they accuse of ``steep price increases'' for VMware's users \citep{reuters2026broadcom}; the F-1 names Broadcom among the company's role models. The portfolio overlaps in video (Vimeo, Brightcove, and StreamYard, a live-streaming and recording platform \citep{silicon2024streamyard}) and in events (Meetup \citep{businesswire2024meetup} and Eventbrite). We found no public record of an agency challenge to any of the company's acquisitions. The F-1's risk factors on antitrust, foreign-investment review and consumer law are not in our extract, so $p^{\rm ch}$, $p^{\rm fdi}$, $T^{\rm fdi}$ and $h^{\rm cp}$ remain scenario parameters.

\subsection{The case: Bending Spoons}\label{sec:acq-case}

Bending Spoons, an Italian software company, describes its playbook as ``acquire, transform and optimize, then reinvest'', essentially unchanged since 2013, and says it reorganizes ``at least with each acquisition'', often into ``a much smaller, more talent-dense organization'' \citep{bendingspoons2026f1}. Table~\ref{tab:acq-case} lists the facts used, each with its source, and Table~\ref{tab:acq-levers} maps its levers to the model's parameters.

\begin{table}[p]
\centering
\caption{The case in figures. F: Form F-1 \citep{bendingspoons2026f1}, whose figures are unchanged in the final prospectus; P: final prospectus \citep{bendingspoons2026prospectus}; K: reports on Form 6-K \citep{bendingspoons2026sixk}; R: reported by the secondary source named, not confirmed in the filings}\label{tab:acq-case}
\scriptsize
\setlength{\tabcolsep}{3pt}
\begin{tabular}{@{}p{3.1cm}p{9.6cm}p{2.9cm}@{}}
\toprule
Item & Value & Source\\
\midrule
Revenue & \acqRevTwentyThree{} (2023), \acqRevTwentyFour{} (2024), \acqRevTwentyFive{} (2025); \acqRevQOneSix{} in Q1 2026 against \acqRevQOneFive & F\\
Margins and interest, 2025 & operating margin \acqOpmTwentyFive; interest expense \acqIntAbsTwentyFive{} (\acqIntTwentyFive{} of revenue), interest paid \acqIntPaidTwentyFive{} (\acqIntPaidShTwentyFive{} of revenue); net income \acqNiTwentyFive; operating cash flow \acqOcfTwentyFive, after interest paid (U.S.\ GAAP); net investing outflow \acqCfiAbsTwentyFive, almost all acquisitions (\acqAcqShareHOneFive{} in the first half); added back in the non-GAAP measures: amortization of acquired intangibles \acqAmzTwentyFive, transaction \acqTrxTwentyFive{} and reorganization \acqReorgTwentyFive{} of revenue (adjusted margin \acqAoimTwentyFive) & F, P\\
Return hurdles and method & \acqHurL{} levered, \acqHurU{} unlevered, for acquisitions completed from 2023 to Q1 2026; \acqTirr{} years of free cash flow plus a terminal value (growth to year 10 and a WACC for recent deals, an EV/EBITDA multiple at year 5 for earlier ones); ``assuming we'll never sell'' & F\\
Debt policy & assumed debt the lower of \acqLtv{} of enterprise value and what projected free cash flow repays in \acqTrep{} years; term loans, expanded after the second quarter by a EUR~500 million loan guaranteed by SACE, EUR~495 million of term loan A, a EUR~490 million increase in revolving commitments (28 July) and add-ons of \acqTlbUsd{} and \acqTlbEur{} to the term loans B due 2031 (2 October), and from 28 October 2026 senior secured notes priced by two subsidiaries, which the company does not issue or guarantee (\acqNotesUsd{} at \acqNotesUsdC{} and \acqNotesEur{} at \acqNotesEurC, due \acqNotesDue), to fund future acquisitions; interest expense of the first half of 2026 at an annual rate plus the notes' coupons, \$\acqRunRateInt{} million a year, is \acqRunRateShQ{} of second-quarter revenue at an annual rate, and \acqRunRateShQHi{} if the new term loans, whose rates the reports do not state, carry the notes' coupon; leverage covenant at most \acqCov, actual \acqLevTwentyFive{} (end-2025), \acqLevQOne{} (end-Q1 2026) and \acqLevQTwo{} (mid-2026, EBITDA with expected cost savings); net debt \acqNetDebt{} (mid-2026) & F, K\\
Capital and equity & \acqDepQOne{} deployed in acquisitions in Q1 2026 against \acqDepTwentyThree{} in 2023; more than \acqTargets{} targets evaluated in 2025; \acqEqRaised{} of primary equity in total, \acqEqShare{} since 2023; dilution \acqDilution{} a year; total assets \acqTotalAssets, liabilities \acqLiabShare{} of assets (31 Mar 2026) & F\\
Core team & \acqSpo{} FTE ``Spooners'' (\acqSpoShare{} of FTE, about \acqFteTotal{} in all; \acqSpoTech{} technical), ``transferred between businesses on short notice''; \acqApps{} applications and \acqHires{} hires in 2025; Remini rewritten in full by about \acqReminiFte{} Spooners after an asset deal & F\\
AI & pull requests written or co-written by AI from under \acqPrLo{} (Q1 2025) to over \acqPrHi{} (end-Q1 2026), about \acqPrSolo{} by AI alone; revenue per Spooner \acqRpsTwentyThree{} (2023), \acqRpsTwentyFive{} (2025), \acqRpsQOne{} (Q1 2026, a quarter and a different denominator); AI helped ``in part'' & F\\
Retention and tenure & portfolio NRR \acqNrrTwentyThree, \acqNrrTwentyFour, \acqNrrTwentyFive, \acqNrrQOne{} (2023 to Q1 2026); by business \acqNrrAol{} (AOL), \acqNrrEvernote{} (Evernote), \acqNrrRemini{} (Remini), \acqNrrStreamyard{} (StreamYard); tenure as in Section~\ref{sec:acq-retention}; \acqOrgTwentyFive{} of new-customer revenue from organic channels (2025) & F\\
Controls and governance & ``We have identified material weaknesses in our internal control over financial reporting'', which integration may leave unremediated or add to; no SOX 404 evaluation yet; Q1 2026 cash flows restated; foreign private issuer, dual class with founders holding class A & F, P\\
IPO & \acqIpoSharesP{} shares at \acqIpoPx, \acqIpoNewP{} of them new, estimated net proceeds to the company \acqIpoNetEst, or \acqIpoNetFull{} if the underwriters' option is exercised in full (P); net proceeds reported as \acqIpoNet{} (K); Berkshire Hathaway \acqBrkOne{} at the IPO and \acqBrkTwo{} later, a \acqBrkStake{} stake (K); \acqIpoTotal{} in all, Nasdaq ticker BSP from 1 July 2026; first day up \acqDayOne; valuation \acqCapOpen{} at the start of trading and \acqCapClose{} at the close; prior private round at \acqPrivRound & P, K; R: \citet{dealroom2026ipo,axios2026ipo,quoteddata2026ipo}\\
Q2 2026 & revenue \acqQTwoRev, up \acqQTwoG; organic growth \acqQTwoOrg{} on the company's definition, unaudited, with estimated pre-acquisition revenue, not presented as a non-GAAP measure; operating margin \acqQTwoM{} (adjusted \acqQTwoAdjM); interest \acqQTwoIntAbs{} (\acqQTwoInt{} of revenue); reorganization \acqQTwoReorgAbs{} (\acqQTwoReorgX); interest paid \acqIntPaidHOneSix{} in the first half, \acqIntPaidRatioHOneSix{} of interest expense & K\\
Workforce actions & Evernote: most staff in the U.S.\ and Chile told of layoffs, operations moved to Europe (July 2023); Mosaic apps: none of \acqMosaic{} staff expected to join; WeTransfer: \acqWetPlanned{} of staff to be laid off (planned, confirmed by the company; over \acqWetStaff{} staff); Brightcove: \acqBcCut{} redundancies (over \acqBcStaff{} staff worldwide); Komoot: up to \acqKomootMax{} of about \acqKomootStaff{} staff reportedly let go, scale not confirmed by the company; Vimeo: a layoff announced on 20 January 2026, no figure confirmed; AOL: \acqAolWarn{} workers at Reston under a WARN notice of 17 February 2026 & R: \citet{appleinsider2023evernote,leach2024mosaic,techcrunch2024wetransfer,tvbeurope2025brightcove,bikeradar2025komoot,petapixel2026vimeo,ffxnow2026aol}\\
Pricing & Harvest: per seat to per seat plus usage; one customer from about \acqHarvestOld{} to \acqHarvestNew{} a month (\acqHarvestRise) & R: \citet{subscriptioninsider2026harvest}, citing BBC News\\
Deal prices & purchase prices AOL \acqPxAol, Eventbrite \acqPxEb, Vimeo \acqPxVim{} (\acqKEb{} to \acqKVim{} times revenue; P); enterprise values Tractive \acqTractive, Airtable \acqAirEv, Miro \acqMiroEv, both completed (K); Vimeo \acqVimeoPx{} a share, \acqVimeoPrem{} above the 60-day average price; Brightcove \acqBcPx{} a share, about \acqBcPrice; Eventbrite \acqEbPx{} a share; AOL with a \acqAolDebt{} debt financing & P, K; R: \citet{vimeo2025release,streamingmedia2024brightcove,marketscreener2026eventbrite,businesswire2025aol}\\
Errors caught & a claim of \acqKucoin{} of 2025 revenue (F-1: \acqRevTwentyFive); a trailing net loss of \acqScoopLoss{} to March 2026 (F-1 tables: net income of \acqTtmNi); Wikipedia's \acqTractiveWiki{} for Tractive (K: \acqTractive) & secondary; see notes\\
\bottomrule
\end{tabular}
\end{table}

\paragraph{What the case is consistent with.} Owner capability: a mobile core team, technology and data presented as one ``Platform'', and a full rewrite of an acquired product by a small team are Assumption~\ref{ass:acq-owner}. Replacement: restructuring with each acquisition and the reported workforce actions are the path of Proposition~\ref{prop:acq-replace}, and reorganization-related expense of \acqReorgTwentyFive{} of revenue in 2025 prices it; the reported cuts span the United States, the Netherlands and Germany, so the case does not show whether labor law shaped where or how fast the owner cut. Permanent capital: never selling, a \acqTirr-year IRR with a terminal value and debt sized to \acqTrep{} years of cash flow are Assumption~\ref{ass:acq-pc}; Proposition~\ref{prop:acq-pc}(b) then requires buying well below the value of a business as transformed; prices of \acqKEb{} to \acqKVim{} times revenue for the three deals in the pro forma statements (Brightcove reportedly near one) allow it only with margins and retention near or above the grid's, which post-deal cash flows would have to show. Retention: business-level NRR of \acqNrrRemini{} to \acqNrrEvernote{} and long tenures fit heterogeneous lock-in and repricing of a locked-in base, as the reported Harvest case suggests; the company reports portfolio retention of \acqNrrTwentyThree{} to \acqNrrTwentyFive. Bought growth and interest fit Proposition~\ref{prop:acquired-score}; the loans and notes raised after the second quarter add to the interest behind its after-interest scores, and the notes' stated use, future acquisitions, shows continued reliance on acquisition capital; and portfolio retention below one with 2025 financing inflows of \acqCffShare{} of revenue fit the treadmill of Proposition~\ref{prop:acq-limits}(a) in its form without new customers; on the company's measure, organic growth of \acqQTwoOrg{} in the second quarter of 2026 means new customers more than offset the decline of the subscription base in that quarter. Controls: material weaknesses alongside rapid integration fit the evidence on growth and restructuring \citep{doyle2007determinants,ashbaughskaife2007discovery}.

\paragraph{What it challenges.} Three conventions of Parts~I and II, and a fourth of measurement. Transformation capability is not a property of the firm: readiness, integration cost, success and horizon in Sections~\ref{sec:tasks} and \ref{sec:ext} should be indexed by owner, and the invoice example's negative value is a statement about a standalone firm. The pace limits of Propositions~\ref{prop:attainable}(c) and \ref{prop:expected} bind in-place change, not replacement, so the attainable frontier is owner-specific. Proposition~\ref{prop:reg-lbo} gives the new owner the target's capability; with a platform the maximum premium gains $\Delta V^{\rm cap}_{\mathsf o}/P_0$ \eqref{eq:acq-pmax}. And growth bought with capital across markets is not on Part~I's frontier of spend in one market, so scores that do not charge it are not comparable (Section~\ref{sec:acquired-growth}). Review of a series of deals adds a constraint that Part~I's competition results, which concern one firm in one market, lack: the owner's cumulative share across deals (Proposition~\ref{prop:acq-limits}(c)).

\paragraph{What it cannot show.} It is one successful case, selected because it reached a listing: survivors overstate a strategy's returns \citep{brown1992survivorship}, roll-up companies that went public delivered poor stock returns \citep{brown2005corporate}, and serial acquirers' returns differ sharply by target type \citep{fuller2002what}. The playbook predates generative AI, so AI's contribution is not identified: AI-written code rose from under \acqPrLo{} to over \acqPrHi{} of pull requests in a year while acquisitions, the team and the denominators of revenue per Spooner all changed, and the company says AI helped ``in part''. Deal-level returns, and so whether returns fall with deal size or for listed targets (Proposition~\ref{prop:acq-limits}(b)), and pre-acquisition retention are not reported. Organic growth is reported only for the second quarter of 2026, on the company's own definition, unaudited and built on estimated pre-acquisition revenue (\acqQTwoOrg, against \acqQTwoG{} reported; Section~\ref{sec:acquired-growth}), and one post-IPO quarter does not establish a trend. Costs that precedent documents for other stakeholders, displaced workers' earnings losses \citep{lagaras2026mas}, reputational harm \citep{love2009character} and a parent's practices lowering quality in acquired units \citep{eliason2020how}, enter the model only through $\delta^G$ and the employment preference of Assumption~\ref{ass:reg-regime}(v); in consumer products, private-equity owners raised prices on existing products by about 1\% and grew through new products \citep{fracassi2022barbarians}, and combined acquirer and target returns are positive on average \citep{andrade2001new}.

\begin{table}[t]
\centering
\caption{The case's levers and the model's parameters}\label{tab:acq-levers}
\scriptsize
\setlength{\tabcolsep}{3pt}
\begin{tabular}{@{}p{3.6cm}p{3.7cm}p{3.4cm}p{4.7cm}@{}}
\toprule
Lever (F-1 wording or report) & Parameter & Result & What the case gives\\
\midrule
``Platform'': people, technology, data & $u_F$, $N_{\rm pl}$, $C^{\rm pool}$, $\varpi^{\mathsf o}$ & Prop.~\ref{prop:acq-owner}(a)--(c) & existence, no costs or reuse shares\\
Spooners moved on short notice & $N^{\rm core}$, $s^{\rm tech}$, $K_{\rm chg}$, $T^{\rm lead}$ & Prop.~\ref{prop:acq-owner}(e), Prop.~\ref{prop:labor-talent} & \acqSpo{} FTE, \acqSpoTech{} technical; at most \acqMaxConcurrent{} Remini-sized teams at once\\
Selective hiring, Talent team & hiring rate, team experience in $\varpi^c_p$ & Prop.~\ref{prop:labor-talent}, \eqref{eq:tasks-success} & \acqHires{} hires from \acqApps{} applications (\acqHireRate); \acqTalent{} FTE (\acqTalentShare)\\
Smaller organization, ``startup mode'' & $\mathsf d^{\rm rpl}$, $t^{\rm rpl}$, $\kappa_{\rm doc}$, $L^{\rm ctx}$ & Prop.~\ref{prop:acq-replace} & reorganization expense \acqReorgTwentyFive{} of revenue (2025)\\
Full rewrite (Remini) & $\mathbb 1^{\rm reb}$, $T^{\rm bld}$, $C^{\rm spec}$ & Prop.~\ref{prop:acq-owner}(b), \ref{prop:acq-replace} & about \acqReminiFte{} FTE; no cost or duration\\
Never sell; \acqTirr-year IRR; hurdles & $T_{\rm pc}$, $T^{\rm IRR}$, $\bar{\mathfrak i}_{\rm U}$, $\bar{\mathfrak i}_{\rm L}$, $r_{\rm tv}$ & Prop.~\ref{prop:acq-pc}(a), (b) & hurdles and method; no $r_{\rm tv}$ in the extract\\
Debt sizing, covenant & $\overline{\rm LTV}$, $\mathrm{ND}^{\rm rep}$, $\overline{\rm lev}^{\rm pc}$ & Prop.~\ref{prop:acq-pc}(c), \ref{prop:reg-dip} & \acqLtv, \acqTrep{} years, \acqCov; notes at \acqNotesUsdC\\
Acquired growth & $g^{\rm acq}$, $k^{\rm acq}$, $c_{\rm fin}$ & Prop.~\ref{prop:acquired-score} & pro forma revenue and prices; organic growth for one quarter (company measure, unaudited)\\
Retention and repricing & $G_i$, $\upsilon_{\rm rp}$, $\mathsf S_{\rm sw}$, $\delta^G$, $h^{\rm cp}$ & Prop.~\ref{prop:acq-retention}, \ref{prop:acq-limits}(d) & NRR levels and tenure; prices only reported\\
Larger, listed targets & $S^{\rm tg}$, $\epsilon_{\rm sz}$, $n^{\rm dl}$, $\mathrm{nd}$ & Prop.~\ref{prop:acq-limits}(a), (b) & prices \acqKEb{} to \acqKVim{} times revenue; no deal returns\\
Overlapping niches; foreign review & $p^{\rm ch}$, $L^{\rm ch}$, $p^{\rm fdi}$, $T^{\rm fdi}$ & Prop.~\ref{prop:acq-limits}(c) & overlaps in video and events; no challenge found\\
AI-written code & $F^A$, $e_{\rm bld}$, $e_{\rm mnt}$, $\Gamma$ & Sections~\ref{sec:tasks}, \ref{sec:labor} & shares of pull requests; effect not identified\\
Material weaknesses & $p_{\rm def}$, $P_{\rm mw}(n)$, $a_{\min}$ & Prop.~\ref{prop:acq-replace}(d), \ref{prop:acct-amin} & existence; no description in the extract\\
Dual class, founders & $\mathfrak w_{\rm dc}$ & Prop.~\ref{prop:reg-hybrid}(a) & structure stated\\
\bottomrule
\end{tabular}
\end{table}
